%% kapa: exam=1 round=27 year=2016 period=1 form=A subject=civil_law subject_ko=민법 seq=12 number=12
%% source: Q-Net 감정평가사 기출문제 (cst00309) artlSeq=5169402 file=extracted/1008/5169402/제27회 감정평가사 1차 문제/제27회 감정평가사 1교시문제(A형).pdf
%% conversion: kordoc 4.24.51 -> tools/gen1.py (auto); review=proofread-sonnet; flags=none
\begin{question}{12}
乙은 대리권 없이 행위능력자인 甲의 임의대리인으로 행세하여 甲 소유의 부동산을 丙에게 매매하는 계약을 체결하였다. 이에 관한 설명으로 옳지 않은 것은? (표현대리는 고려하지 않으며, 다툼이 있으면 판례에 따름)
\begin{choices}
  \item 乙이 위 계약에 따라 丙에게 소유권이전등기를 해준 경우, 甲은 丙 명의 등기의 말소를 청구할 수 있다.
  \item 乙이 위 계약 당시 제한능력자인 경우, 乙은 丙에게 계약의 이행 또는 손해배상책임을 지지 않는다.
  \item 甲이 乙의 무권대리행위를 알면서도 丙에게 매매대금을 청구하여 전부를 수령하였다면, 특별한 사정이 없는 한, 위 계약을 추인한 것으로 볼 수 있다.
  \item 甲이 乙에 대하여 추인을 하였다면 丙이 그 추인 사실을 몰랐더라도 위 계약을 철회할 수 없다.
  \item 甲의 유효한 추인이 있으면, 특별한 사정이 없는 한, 乙의 행위는 계약 시에 소급하여 甲에게 효력이 있다.
\end{choices}
\end{question}
\answer{4}
